\documentclass[11pt]{article}

\usepackage{amsmath, amssymb}
\usepackage{graphicx}
\usepackage{geometry}
\usepackage{setspace}
\usepackage{titlesec}

\geometry{margin=1in}

\titleformat{\section}{\large\bfseries}{\thesection}{1em}{}
\titleformat{\subsection}{\normalsize\bfseries}{\thesubsection}{1em}{}

\title{
CSE400 -- Fundamentals of Probability in Computing\\
\large Lecture 12 Scribe: Transformation of Random Variables and Functions of Two Random Variables
}

\author{
Dhaval Patel, PhD\\
Associate Professor, Computer Science and Engineering\\
SEAS -- Ahmedabad University
}

\date{}

\begin{document}

\maketitle

\onehalfspacing

\section{Lecture Overview}

This lecture studies how the probability distribution of a random variable changes when a transformation is applied. The lecture focuses on three major topics:

\begin{enumerate}
    \item Transformation of Random Variables
    \item Functions of Two Random Variables
    \item Illustrative Example: Derivation of the distribution of $Z = X + Y$
\end{enumerate}

The objective is to determine the probability distribution of a new random variable when it is defined as a function of one or more existing random variables.

\section{Transformation of a Random Variable}

\subsection{Problem Setup}

Suppose a continuous random variable $X$ has known probability density function (PDF)

\[
f_X(x)
\]

and cumulative distribution function (CDF)

\[
F_X(x)
\]

Now define a new random variable

\[
Y = g(X)
\]

The goal is to determine the distribution of $Y$, namely

\[
F_Y(y) \quad \text{and} \quad f_Y(y)
\]

\section{Derivation Using the CDF Method}

\subsection{Step 1: Definition of CDF}

The cumulative distribution function of $Y$ is

\[
F_Y(y) = P(Y \le y)
\]

Substituting $Y = g(X)$ gives

\[
F_Y(y) = P(g(X) \le y)
\]

\subsection{Step 2: Use the inverse transformation}

If the function $g(\cdot)$ is monotonic, the inverse function exists. Therefore,

\[
g(X) \le y
\]

is equivalent to

\[
X \le g^{-1}(y)
\]

Hence,

\[
F_Y(y) = P(X \le g^{-1}(y))
\]

\[
F_Y(y) = F_X(g^{-1}(y))
\]

\section{Deriving the PDF}

To obtain the PDF of $Y$, differentiate the CDF.

\[
f_Y(y) = \frac{d}{dy} F_Y(y)
\]

Substituting the previous result,

\[
f_Y(y) = \frac{d}{dy} F_X(g^{-1}(y))
\]

Using the chain rule,

\[
f_Y(y) =
f_X(g^{-1}(y))
\frac{d}{dy} g^{-1}(y)
\]

Thus the transformation formula becomes

\[
f_Y(y) =
f_X(x)
\left|
\frac{dx}{dy}
\right|
\quad \text{where } x = g^{-1}(y)
\]

\section{Example: Transformation of a Uniform Random Variable}

\subsection{Given}

\[
X \sim U(-1,1)
\]

Thus the PDF of $X$ is

\[
f_X(x) =
\begin{cases}
\frac{1}{2}, & -1 < x < 1 \\
0, & \text{otherwise}
\end{cases}
\]

Define the transformation

\[
Y = \sin\left(\frac{\pi X}{2}\right)
\]

Find $f_Y(y)$.

\subsection{Step 1: Invert the transformation}

\[
Y = \sin\left(\frac{\pi X}{2}\right)
\]

Taking the inverse sine,

\[
X = \frac{2}{\pi}\sin^{-1}(y)
\]

\subsection{Step 2: Compute derivative}

\[
\frac{dx}{dy} =
\frac{2}{\pi}\frac{1}{\sqrt{1-y^2}}
\]

\subsection{Step 3: Apply transformation formula}

\[
f_Y(y) =
f_X(x)\left|\frac{dx}{dy}\right|
\]

Substitute $f_X(x) = \frac{1}{2}$

\[
f_Y(y) =
\frac{1}{2}
\cdot
\frac{2}{\pi}
\frac{1}{\sqrt{1-y^2}}
\]

\subsection{Final Result}

\[
f_Y(y) =
\frac{1}{\pi\sqrt{1-y^2}}, \quad -1 < y < 1
\]

\section{Function of Two Random Variables}

Now consider a transformation involving two random variables:

\[
Z = X + Y
\]

The goal is to determine the probability density function of $Z$.

\section{CDF of the Transformed Variable}

Start with the definition

\[
F_Z(z) = P(Z \le z)
\]

Substituting $Z = X + Y$

\[
F_Z(z) = P(X + Y \le z)
\]

This probability corresponds to the region

\[
x + y \le z
\]

in the $(x,y)$ plane.

Thus,

\[
F_Z(z) =
\iint_{x+y\le z} f_{XY}(x,y)\,dx\,dy
\]

\section{Differentiation Using Leibniz Rule}

To obtain the PDF,

\[
f_Z(z) = \frac{d}{dz}F_Z(z)
\]

Using the Leibniz rule for differentiation of integrals,

\[
f_Z(z) =
\int_{-\infty}^{\infty}
f_{XY}(x,z-x)\,dx
\]

\section{Independent Random Variables (Convolution)}

If $X$ and $Y$ are independent,

\[
f_{XY}(x,y) = f_X(x)f_Y(y)
\]

Therefore,

\[
f_Z(z) =
\int_{-\infty}^{\infty}
f_X(x)f_Y(z-x)\,dx
\]

This integral is called the convolution integral.

\section{Example: Sum of Two Normal Random Variables}

Let

\[
X \sim N(0,1)
\]

\[
Y \sim N(0,1)
\]

with independence.

Applying the convolution formula yields

\[
f_Z(z) =
\frac{1}{\sqrt{4\pi}}e^{-z^2/4}
\]

Thus,

\[
Z \sim N(0,2)
\]

\section{Example: Sum of Exponential Random Variables}

Let

\[
X,Y \sim \text{Exponential}(\lambda)
\]

with

\[
f_X(x) = \lambda e^{-\lambda x}, \quad x>0
\]

\[
f_Y(y) = \lambda e^{-\lambda y}, \quad y>0
\]

Using convolution,

\[
f_Z(z) =
\int_0^z
\lambda e^{-\lambda x}
\lambda e^{-\lambda (z-x)} dx
\]

which simplifies to

\[
f_Z(z) =
\lambda^2 z e^{-\lambda z}, \quad z>0
\]

\section{Additional Transformations}

Other transformations discussed include

\[
Z = X - Y
\]

\[
Z = \frac{X}{Y}
\]

and radial transformations

\[
R = \sqrt{X^2 + Y^2}
\]

These illustrate how functions of random variables can be used to derive new probability distributions.

\section{Key Results for Examination}

\subsection{Single Variable Transformation}

\[
f_Y(y) =
f_X(g^{-1}(y))
\left|
\frac{d}{dy}g^{-1}(y)
\right|
\]

\subsection{Distribution of the Sum}

\[
f_Z(z) =
\int f_{XY}(x,z-x)dx
\]

\subsection{Independent Case (Convolution)}

\[
f_Z(z) =
\int f_X(x)f_Y(z-x)dx
\]

\subsection{Important Results}

\[
X,Y \sim N(0,1)
\Rightarrow
Z = X+Y \sim N(0,2)
\]

\[
X,Y \sim \text{Exponential}(\lambda)
\Rightarrow
f_Z(z) = \lambda^2 z e^{-\lambda z}
\]

\end{document}
